% Options for packages loaded elsewhere
\PassOptionsToPackage{unicode}{hyperref}
\PassOptionsToPackage{hyphens}{url}
\documentclass[
]{article}
\usepackage{xcolor}
\usepackage[margin=1in]{geometry}
\usepackage{amsmath,amssymb}
\setcounter{secnumdepth}{-\maxdimen} % remove section numbering
\usepackage{iftex}
\ifPDFTeX
  \usepackage[T1]{fontenc}
  \usepackage[utf8]{inputenc}
  \usepackage{textcomp} % provide euro and other symbols
\else % if luatex or xetex
  \usepackage{unicode-math} % this also loads fontspec
  \defaultfontfeatures{Scale=MatchLowercase}
  \defaultfontfeatures[\rmfamily]{Ligatures=TeX,Scale=1}
\fi
\usepackage{lmodern}
\ifPDFTeX\else
  % xetex/luatex font selection
\fi
% Use upquote if available, for straight quotes in verbatim environments
\IfFileExists{upquote.sty}{\usepackage{upquote}}{}
\IfFileExists{microtype.sty}{% use microtype if available
  \usepackage[]{microtype}
  \UseMicrotypeSet[protrusion]{basicmath} % disable protrusion for tt fonts
}{}
\makeatletter
\@ifundefined{KOMAClassName}{% if non-KOMA class
  \IfFileExists{parskip.sty}{%
    \usepackage{parskip}
  }{% else
    \setlength{\parindent}{0pt}
    \setlength{\parskip}{6pt plus 2pt minus 1pt}}
}{% if KOMA class
  \KOMAoptions{parskip=half}}
\makeatother
\usepackage{color}
\usepackage{fancyvrb}
\newcommand{\VerbBar}{|}
\newcommand{\VERB}{\Verb[commandchars=\\\{\}]}
\DefineVerbatimEnvironment{Highlighting}{Verbatim}{commandchars=\\\{\}}
% Add ',fontsize=\small' for more characters per line
\usepackage{framed}
\definecolor{shadecolor}{RGB}{248,248,248}
\newenvironment{Shaded}{\begin{snugshade}}{\end{snugshade}}
\newcommand{\AlertTok}[1]{\textcolor[rgb]{0.94,0.16,0.16}{#1}}
\newcommand{\AnnotationTok}[1]{\textcolor[rgb]{0.56,0.35,0.01}{\textbf{\textit{#1}}}}
\newcommand{\AttributeTok}[1]{\textcolor[rgb]{0.13,0.29,0.53}{#1}}
\newcommand{\BaseNTok}[1]{\textcolor[rgb]{0.00,0.00,0.81}{#1}}
\newcommand{\BuiltInTok}[1]{#1}
\newcommand{\CharTok}[1]{\textcolor[rgb]{0.31,0.60,0.02}{#1}}
\newcommand{\CommentTok}[1]{\textcolor[rgb]{0.56,0.35,0.01}{\textit{#1}}}
\newcommand{\CommentVarTok}[1]{\textcolor[rgb]{0.56,0.35,0.01}{\textbf{\textit{#1}}}}
\newcommand{\ConstantTok}[1]{\textcolor[rgb]{0.56,0.35,0.01}{#1}}
\newcommand{\ControlFlowTok}[1]{\textcolor[rgb]{0.13,0.29,0.53}{\textbf{#1}}}
\newcommand{\DataTypeTok}[1]{\textcolor[rgb]{0.13,0.29,0.53}{#1}}
\newcommand{\DecValTok}[1]{\textcolor[rgb]{0.00,0.00,0.81}{#1}}
\newcommand{\DocumentationTok}[1]{\textcolor[rgb]{0.56,0.35,0.01}{\textbf{\textit{#1}}}}
\newcommand{\ErrorTok}[1]{\textcolor[rgb]{0.64,0.00,0.00}{\textbf{#1}}}
\newcommand{\ExtensionTok}[1]{#1}
\newcommand{\FloatTok}[1]{\textcolor[rgb]{0.00,0.00,0.81}{#1}}
\newcommand{\FunctionTok}[1]{\textcolor[rgb]{0.13,0.29,0.53}{\textbf{#1}}}
\newcommand{\ImportTok}[1]{#1}
\newcommand{\InformationTok}[1]{\textcolor[rgb]{0.56,0.35,0.01}{\textbf{\textit{#1}}}}
\newcommand{\KeywordTok}[1]{\textcolor[rgb]{0.13,0.29,0.53}{\textbf{#1}}}
\newcommand{\NormalTok}[1]{#1}
\newcommand{\OperatorTok}[1]{\textcolor[rgb]{0.81,0.36,0.00}{\textbf{#1}}}
\newcommand{\OtherTok}[1]{\textcolor[rgb]{0.56,0.35,0.01}{#1}}
\newcommand{\PreprocessorTok}[1]{\textcolor[rgb]{0.56,0.35,0.01}{\textit{#1}}}
\newcommand{\RegionMarkerTok}[1]{#1}
\newcommand{\SpecialCharTok}[1]{\textcolor[rgb]{0.81,0.36,0.00}{\textbf{#1}}}
\newcommand{\SpecialStringTok}[1]{\textcolor[rgb]{0.31,0.60,0.02}{#1}}
\newcommand{\StringTok}[1]{\textcolor[rgb]{0.31,0.60,0.02}{#1}}
\newcommand{\VariableTok}[1]{\textcolor[rgb]{0.00,0.00,0.00}{#1}}
\newcommand{\VerbatimStringTok}[1]{\textcolor[rgb]{0.31,0.60,0.02}{#1}}
\newcommand{\WarningTok}[1]{\textcolor[rgb]{0.56,0.35,0.01}{\textbf{\textit{#1}}}}
\usepackage{graphicx}
\makeatletter
\newsavebox\pandoc@box
\newcommand*\pandocbounded[1]{% scales image to fit in text height/width
  \sbox\pandoc@box{#1}%
  \Gscale@div\@tempa{\textheight}{\dimexpr\ht\pandoc@box+\dp\pandoc@box\relax}%
  \Gscale@div\@tempb{\linewidth}{\wd\pandoc@box}%
  \ifdim\@tempb\p@<\@tempa\p@\let\@tempa\@tempb\fi% select the smaller of both
  \ifdim\@tempa\p@<\p@\scalebox{\@tempa}{\usebox\pandoc@box}%
  \else\usebox{\pandoc@box}%
  \fi%
}
% Set default figure placement to htbp
\def\fps@figure{htbp}
\makeatother
\setlength{\emergencystretch}{3em} % prevent overfull lines
\providecommand{\tightlist}{%
  \setlength{\itemsep}{0pt}\setlength{\parskip}{0pt}}
\usepackage{bookmark}
\IfFileExists{xurl.sty}{\usepackage{xurl}}{} % add URL line breaks if available
\urlstyle{same}
\hypersetup{
  pdftitle={rstanarm},
  hidelinks,
  pdfcreator={LaTeX via pandoc}}

\title{rstanarm}
\author{}
\date{\vspace{-2.5em}2026-04-17}

\begin{document}
\maketitle

\subsection{Introduction}\label{introduction}

\texttt{rstanarm} is a sister package to \texttt{brms} from the Stan
team. Where \texttt{brms} translates each formula into a fresh Stan
program and compiles it on the fly, \texttt{rstanarm} ships with a fixed
catalogue of pre-compiled models --- \texttt{stan\_lm},
\texttt{stan\_glm}, \texttt{stan\_glmer}, \texttt{stan\_polr},
\texttt{stan\_betareg}, and a handful of others --- that match the most
common R regression interfaces. The trade-off is straightforward: no
compilation latency at fit time, weakly informative defaults curated by
experts, and a learning curve identical to base-R modelling functions,
in exchange for less flexibility when the model deviates from the
standard catalogue.

\subsection{Prerequisites}\label{prerequisites}

A working knowledge of generalised linear models, an idea of what a
prior and a posterior are, and comfort with R's \texttt{lm()} /
\texttt{glm()} / \texttt{lmer()} syntax. No C++ toolchain is needed
beyond what installing the package itself requires, since the Stan
models are pre-compiled.

\subsection{Theory}\label{theory}

Each \texttt{stan\_*} function is a thin R wrapper around a Stan program
that has been compiled when the package was built. The program already
includes weakly informative default priors centred at zero with scales
determined by the standard deviation of each predictor. Override
defaults through the \texttt{prior}, \texttt{prior\_intercept},
\texttt{prior\_aux}, and \texttt{prior\_covariance} arguments, each
accepting one of \texttt{normal()}, \texttt{student\_t()},
\texttt{cauchy()}, \texttt{laplace()}, \texttt{exponential()}, or
\texttt{decov()}. Because the priors are scaled to the predictors, the
same numerical scale (e.g.~\texttt{normal(0,\ 2.5)}) is sensible across
very different predictor ranges.

\subsection{Assumptions}\label{assumptions}

The assumptions of the underlying GLM or GLMM apply: linearity in the
linear predictor, independence after grouping, the appropriate link /
variance relationship for the family. Convergence (\(\hat R < 1.01\), no
divergent transitions) is a precondition for trusting the posterior
summary that \texttt{rstanarm} returns.

\subsection{R Implementation}\label{r-implementation}

\begin{Shaded}
\begin{Highlighting}[]
\FunctionTok{library}\NormalTok{(rstanarm)}

\FunctionTok{set.seed}\NormalTok{(}\DecValTok{2026}\NormalTok{)}
\NormalTok{d }\OtherTok{\textless{}{-}} \FunctionTok{data.frame}\NormalTok{(}\AttributeTok{y =} \FunctionTok{rnorm}\NormalTok{(}\DecValTok{100}\NormalTok{), }\AttributeTok{x =} \FunctionTok{rnorm}\NormalTok{(}\DecValTok{100}\NormalTok{))}

\CommentTok{\# Same syntax as lm() but Bayesian}
\NormalTok{fit }\OtherTok{\textless{}{-}} \FunctionTok{stan\_glm}\NormalTok{(y }\SpecialCharTok{\textasciitilde{}}\NormalTok{ x, }\AttributeTok{data =}\NormalTok{ d,}
                \AttributeTok{prior =} \FunctionTok{normal}\NormalTok{(}\DecValTok{0}\NormalTok{, }\FloatTok{2.5}\NormalTok{),}
                \AttributeTok{prior\_intercept =} \FunctionTok{normal}\NormalTok{(}\DecValTok{0}\NormalTok{, }\DecValTok{5}\NormalTok{),}
                \AttributeTok{prior\_aux =} \FunctionTok{exponential}\NormalTok{(}\DecValTok{1}\NormalTok{),}
                \AttributeTok{chains =} \DecValTok{4}\NormalTok{, }\AttributeTok{iter =} \DecValTok{2000}\NormalTok{, }\AttributeTok{refresh =} \DecValTok{0}\NormalTok{)}
\FunctionTok{summary}\NormalTok{(fit)}
\FunctionTok{posterior\_interval}\NormalTok{(fit, }\AttributeTok{prob =} \FloatTok{0.95}\NormalTok{)}
\end{Highlighting}
\end{Shaded}

\subsection{Output \& Results}\label{output-results}

\texttt{summary(fit)} lists the posterior mean, standard deviation, 10 /
50 / 90 \% quantiles, mean posterior predictive checks (PPP),
\(\hat R\), and effective sample size for each parameter.
\texttt{posterior\_interval()} returns the requested credible interval.
Because the same Stan program is reused, fits typically take seconds
rather than minutes.

\subsection{Interpretation}\label{interpretation}

A reporting sentence: ``A Bayesian linear regression
(\texttt{stan\_glm}, weakly informative priors, four chains × 2,000
iterations) estimated the slope at 0.47 (95 \% CrI 0.27 to 0.68); the
posterior overlaps the OLS point estimate but quantifies the residual
uncertainty rather than just standard errors.'' Compared with
frequentist output, the credible interval has the direct probability
interpretation that students often (incorrectly) assign to confidence
intervals.

\subsection{Practical Tips}\label{practical-tips}

\begin{itemize}
\tightlist
\item
  Use \texttt{prior\_summary(fit)} to verify exactly which priors were
  applied --- the auto-scaling can surprise you when predictors are on
  unusual scales.
\item
  For mixed models, \texttt{stan\_glmer()} mirrors
  \texttt{lme4::glmer()} syntax including correlated random slopes.
\item
  \texttt{pp\_check(fit)} provides graphical posterior predictive
  checks; \texttt{bayes\_R2()} returns the Bayesian counterpart of
  \(R^2\) as a posterior distribution.
\item
  \texttt{loo(fit)} and \texttt{waic(fit)} attach information criteria
  with Pareto-\(k\) diagnostics, ready for \texttt{loo\_compare()}
  between models.
\item
  When the model you need is not in the catalogue (for example, custom
  likelihoods, censored data, ordinal-with-arbitrary-link), reach for
  \texttt{brms} or raw Stan; \texttt{rstanarm} is deliberately scoped.
\item
  For very small datasets, prefer \texttt{student\_t(df\ =\ 3)} priors
  over \texttt{normal()} to reduce posterior sensitivity to a single
  influential observation.
\end{itemize}

\subsection{R Packages Used}\label{r-packages-used}

\texttt{rstanarm} for the pre-compiled Stan models, \texttt{bayesplot}
for diagnostic plots underlying \texttt{pp\_check()}, \texttt{loo} for
cross-validation, and \texttt{posterior} for working with the underlying
draws.

\subsection{For Reviewers}\label{for-reviewers}

\textbf{What to look for in a paper using this method.}

\begin{itemize}
\tightlist
\item
  \textbf{Common misapplications.}
\item
  \textbf{Diagnostics that should be reported but often aren't.}
\item
  \textbf{Red flags in tables and figures.}
\item
  \textbf{What to verify.}
\item
  \textbf{What an adequate Methods paragraph must contain.}
\end{itemize}

\subsection{See also --- labs in this
chapter}\label{see-also-labs-in-this-chapter}

\begin{itemize}
\tightlist
\item
  \href{labs/lab-bayesian-thinking-control-vs-decision-error.Rmd}{Bayesian
  thinking; α control vs decision error}
\item
  \href{labs/lab-brms-stan-loo-hierarchical-models.Rmd}{brms / Stan,
  LOO, hierarchical models}
\end{itemize}

\end{document}
